\section{Justificación}

\begin{frame}{Justificación}
	\begin{center}
		\begin{tikzpicture}
			\node[cajaobjetivo, text width=0.85\textwidth]
			{
				\textbf{Reducción de complejidad} \\
				Simular sistemas complejos sin realizar ajustes manuales en la función de transición, un proceso que puede resultar computacionalmente inviable.
			};
		\end{tikzpicture}
	\end{center}
	
	\begin{center}
		\begin{tikzpicture}
			\node[cajaobjetivo, text width=0.85\textwidth]
			{
				\textbf{Nuevo alcance} \\
				Combinar herramientas de inteligencia artificial y simulación de sistemas complejos para ampliar el campo de estudio de los autómatas celulares.
			};
		\end{tikzpicture}
	\end{center}
\end{frame}

\begin{frame}{Justificación}
	\begin{center}
		\begin{tikzpicture}
			\node[cajaobjetivo, text width=0.85\textwidth]
			{
				\textbf{Innovación} \\
				Reemplazar funciones de transición estáticas por modelos ajustables mediante técnicas de aprendizaje automático, permitiendo la adaptación del sistema a cambios en el entorno.
			};
		\end{tikzpicture}
	\end{center}
	
	\begin{center}
		\begin{tikzpicture}
			\node[cajaobjetivo, text width=0.85\textwidth]
			{
				\textbf{Tecnología} \\
				Las redes neuronales, especialmente convolucionales y recurrentes, son efectivas para aproximar funciones no lineales y capturar relaciones complejas en espacio y tiempo.
			};
		\end{tikzpicture}
	\end{center}
\end{frame}

\begin{frame}{Justificación}
	\begin{center}
		\begin{tikzpicture}
			\node[cajaobjetivo, text width=0.85\textwidth]
			{
				\textbf{Formación integral} \\
				Profundizar en simulación, modelado matemático y redes neuronales, integrando conocimientos de ciencias de la computación, vida artificial y matemáticas.
			};
		\end{tikzpicture}
	\end{center}
	
	\begin{center}
		\begin{tikzpicture}
			\node[cajaobjetivo, text width=0.85\textwidth]
			{
				\textbf{Aporte al estado del arte} \\
				Fomentar futuras investigaciones mediante la creación de un entorno de código abierto y contribuir al avance en inteligencia artificial y sistemas complejos.
			};
		\end{tikzpicture}
	\end{center}
\end{frame}
